\documentclass[journal, twocolumn]{IEEEtran}

\usepackage{cite}
\usepackage{amsmath,amssymb,amsfonts}
\usepackage{algorithmic}
\usepackage{graphicx}
\usepackage{textcomp}
\usepackage{xcolor}
\usepackage{booktabs}
\usepackage{multirow}

\begin{document}

\title{Revealing Hidden Topological Heterogeneity in Wearable Sensor Data through Dynamic Baseline Anchoring}

\author{
    \IEEEauthorblockN{Kotaro Furukawa, \textit{Student Member, IEEE}}\\
    \IEEEauthorblockA{Frontier School of Transdisciplinary Sciences, Kanazawa University, Ishikawa, Japan\\
    Email: f.kotaro.0530@gmail.com}
}

\maketitle

\begin{abstract}
Continuous physiological sensing is foundational to affective computing and ambulatory monitoring. Standard processing pipelines typically employ generic temporal normalization (e.g., rolling averages) to mitigate hardware-induced drift. We hypothesize that these standard normalization methods artificially enforce population-level topological homogeneity, masking individual-specific covariance structures inherent in biological responses. In this letter, we test this hypothesis by comparing conventional windowed normalizations against a high-resolution (100 Hz), Dynamic Baseline Anchoring (DBA) methodology based on individual Mahalanobis distance. Applying this to 45 subjects across the WESAD and CASE datasets, we demonstrate that DBA reveals a statistically significant bimodal distribution (Hartigan's dip test, $p < 0.05$) in the geometric overlap of physiological stress topologies. While physiological state separability (mean F1: 0.86\textendash0.89) is strictly maintained across the population, the topological overlap between resting and active task responses exhibits extreme dichotomy, ranging from near 0\% to nearly 100\% dependent entirely on the subject. These findings highlight that sensor calibration design fundamentally shapes the observable data topology, potentially introducing structural distortions under conventional normalization pipelines.
\end{abstract}

\begin{IEEEkeywords}
Wearable Sensors, Sensor Calibration, Data Normalization, Physiological Topology, Electrodermal Activity (EDA), Heterogeneity.
\end{IEEEkeywords}

\section{Introduction}
\IEEEPARstart{W}{earable} physiological sensors, integrating real-time modalities such as electrocardiogram (ECG) and electrodermal activity (EDA), assert significant potential for continuous objective state monitoring \cite{picard1997}. However, deploying these sensors in uncontrolled environments introduces substantial signal variability originating from both dynamic hardware contact and vast intra-subject autonomic variance \cite{wilhelm2006}. 

Conventional multi-sensor systems address this variability by universally processing raw outputs utilizing global standardization (Z-score) or localized rolling-window averaging \cite{schmidt2018wesad}. While these normalization techniques successfully bound continuous data conceptually, we hypothesize that coarse temporal windowing (e.g., 60-second intervals) artificially homogenizes instantaneous physiological events, blurring the true geometric topology of the sensor output.

Standard normalization essentially operates on a potentially limiting assumption: the presumption of uniform measurement topology across populations. Rather than attributing variability solely to biological differences, this work frames the observed heterogeneity as a consequence of calibration design within sensor processing pipelines. To explore this preprocessing-induced distortion, we investigate how differing calibration methodologies impact the observed multidimensional signal geometry.

The core contribution of this letter is the quantitative discovery that applying an individually covariant calibration exposes latent topological heterogeneity otherwise suppressed as a calibration artifact by conventional normalization. We explore a Dynamic Baseline Anchoring (DBA) technique that instantaneously aligns multiple sensor streams via the Mahalanobis distance relative to each user's unique resting covariance manifold. By statistically evaluating multi-dimensional signal overlap across heterogeneous datasets, we demonstrate that sensor representational topologies exhibit robust, bimodal variance, fundamentally questioning the deployment of universal static normalization models in precision systems.

\section{Methodology: Sensor Validation Pipeline}

\subsection{Instantaneous Multimodal Extraction}
To analyze specific physiological topologies without temporal smoothing artifacts, we synchronized all modalities into a unified 100 Hz multidimensional feature space.

\subsubsection{Electrodermal Activity (EDA) Deconvolution}
Standard continuous EDA is highly susceptible to massive sweat accumulation drift. Utilizing continuous-time convex optimization (cvxEDA) \cite{cvxeda2015}, the signal is separated into a baseline tonic drift ($EDA_{Tonic}$) and instantaneous sudomotor phasic bursts ($EDA_{Phasic}$). While susceptible to varying skin impedance, extracting these instantaneous components mitigates the reliance on generic moving-window detrending.

\subsubsection{Wavelet-Based Heart Rate Extraction}
We bypass standard windowed Fast Fourier Transforms (FFT) by applying a Continuous Wavelet Transform (CWT) equipped with an analytic Morlet wavelet to the R-R interval time series. This method outputs highly reactive instantaneous vectors for standard frequency bands ($HRV_{LF}$ and $HRV_{HF}$) at 100 Hz, avoiding artificial latency.

\subsection{Dynamic Baseline Anchoring (DBA)}
Extraction generates an instantaneous vector $X_t \in \mathbb{R}^4$. Generic Euclidean scaling applied to this matrix frequently fails, as dimensions with high hardware variance physically dominate the representation space. To establish a scale-invariant structure, DBA anchors data utilizing the covariance matrix.

A 20-minute controlled resting block dictates the specific multivariate noise topology of the user's sensor deployment: $\mathcal{N}(\mu_{base}, \Sigma_{base})$. Subsequent incoming vectors $X_t$ are evaluated continuously using the Mahalanobis distance ($D_M$):
\begin{equation}
    D_M(X_t) = \sqrt{(X_t - \mu_{base})^T \Sigma_{base}^{-1} (X_t - \mu_{base})}
\end{equation}
The inversion matrix $\Sigma_{base}^{-1}$ normalizes the multi-sensor stream based explicitly on the measured resting covariances, effectively neutralizing systemic multidimensional drift artifacts without enforcing arbitrary global window sizes.

\section{Evaluation and Statistical Validation}

\subsection{Baseline Comparisons and Datasets}
To test the hypothesis of artificial homogenization, we analyzed 45 subjects across two separate continuous datasets: WESAD (15 subjects) \cite{schmidt2018wesad} and CASE (30 subjects). We quantitatively compared the 4D topological structures produced by three competing normalizations:
\begin{enumerate}
    \item \textbf{Global Z-Score}: Universal distribution scaling.
    \item \textbf{Rolling Z-Score}: 60-second local moving average.
    \item \textbf{Mahalanobis DBA}: The proposed individually anchored covariant topology.
\end{enumerate}

\subsection{Validating Theoretical Topologies}
To evaluate the impact of calibration, we utilized two metrics:
\begin{itemize}
    \item \textbf{Separability (F1-Score)}: Utilizing logistic regression (Stress vs. Baseline), we measured whether the calibration destroyed underlying biological distinction. 
    \item \textbf{Geometric Overlap ($\Omega$)}: $\Omega$ is computed as the proportion of active task samples whose Mahalanobis distance satisfies $D_M \le \tau$. The threshold $\tau$ (equivalent to $1.0\sigma$) was selected as a natural boundary of the baseline covariance ellipsoid; robustness analysis across $\sigma \in [0.5, 2.0]$ confirmed that the subsequent bimodal distribution persists.
\end{itemize}

\section{Empirical Discoveries}

\subsection{Preservation of Signal Separability}
Table \ref{tab:comparison} details the mean sensor metrics. Critically, conventional Rolling Z-score methods successfully increased topological overlap ($\Omega$), but at the massive cost of destroying physiological signal separability (F1 crashed from 0.86 to 0.61 in WESAD). This confirms our hypothesis that macroscopic temporal windowing artificially homogenizes and crushes acute responses. Conversely, the DBA framework generated massive overlap increases while maintaining peak analytical Separability ($\approx 0.88$).

\begin{table}[tb]
\caption{Mean Topologies Impacted by Preprocessing Methods}
\label{tab:comparison}
\centering
\begin{tabular}{lccc}
\toprule
Dataset & Preprocessing & Mean Overlap ($\Omega$) & Separability (F1) \\
\midrule
\multirow{3}{*}{\textbf{WESAD}} 
& Global Z-Score & 37.6\% & 0.86 \\
& Rolling Z-Score & 65.2\% & 0.61 \\
& Mahalanobis DBA & \textbf{74.9}\% & \textbf{0.86} \\
\midrule
\multirow{2}{*}{\textbf{CASE}}
& Global Z-Score & 37.1\% & 0.89 \\
& Mahalanobis DBA & \textbf{90.0}\% & \textbf{0.89} \\
\bottomrule
\end{tabular}
\end{table}

This result reveals a critical inconsistency: despite high classification separability, the underlying geometric topology can either completely diverge or fully overlap with baseline structures, depending entirely on the individual physiological covariance structure.

\subsection{Exposing Latent Topological Heterogeneity}
While mean overlaps appear indicative of general sensor drift reduction, decomposing the DBA output exposed extreme underlying structures. A Hartigan’s dip test on the individual topological overlap distribution confirmed a massive, statistically significant deviation from unimodality ($p < 0.05$). 

Figure \ref{fig:method_comparison} demonstrates this bimodal polarization. The distribution exhibits clear polarization with peaks near 0\% and 100\%. Gaussian mixture modeling further supports a two-component structural divide. While global and rolling normalizations force subject variance into a generic mediocre distribution ($\approx 30-60\%$), DBA effectively unlocks the underlying sensor data structure. Some subjects exhibited $\Omega \approx 0\%$, generating perfectly distinct topographical separations between resting and stress responses. Paradoxically, numerous other subjects produced $\Omega \approx 100\%$, indicating their massive, fully separable active stress topology geometrically sat flawlessly identically atop their basal resting manifold. An ANOVA tracking intra-subject variance confirmed that the DBA method significantly uncovers latent inter-subject heterogeneity ($p < 0.05$).

\begin{figure}[tb]
    \centering
    \includegraphics[width=0.45\textwidth]{../results/figures/fig2_method_comparison.png}
    \caption{Distribution overlap ($\Omega$) across calibration methodologies. The DBA framework exposes extreme inter-subject bimodal topological differences that are obscured by conventional rolling-window normalization.}
    \label{fig:method_comparison}
\end{figure}

\section{Conclusion and Paradigm Implications}
By utilizing a high-resolution, individually covariant calibration method (DBA), this letter revealed that continuous wearable sensor streams exhibit wildly bimodal geometric topologies beneath generic baseline drift. This study demonstrates that calibration methodology directly governs the observed structure of multimodal sensor data. Rather than sensor representations being uniformly invariant, we discovered that identical tasks produced vastly discordant signal geometries across different sensor representations. We suggest that affective somatic data cannot be mathematically treated as a stationary or uniform mapping, but rather must acknowledge profound, biologically unique topological deviations that are artificially suppressed as preprocessing-induced distortions. 

This study is currently limited to two public datasets and specific feature representations; further validation across additional sensing modalities is required. Therefore, sensor systems must incorporate calibration strategies that preserve individual covariance structures rather than enforcing uniform normalization.

\bibliographystyle{IEEEtran}
\begin{thebibliography}{9}

\bibitem{picard1997}
R. W. Picard, \textit{Affective Computing}. MIT Press, 1997.

\bibitem{wilhelm2006}
P. Wilhelm and M. Schoebi, ``Assessing mood in daily life,'' \textit{European Journal of Psychological Assessment}, vol. 23, no. 4, pp. 258--267, 2007.

\bibitem{barrett2017}
L. F. Barrett, \textit{How emotions are made: The secret life of the brain}. Houghton Mifflin Harcourt, 2017.

\bibitem{schmidt2018wesad}
P. Schmidt, A. Reiss, R. Duerichen, C. Marberger, and K. Van Laerhoven, ``Introducing WESAD, a multimodal dataset for wearable stress and affect detection,'' in \textit{Proceedings of the 20th ACM International Conference on Multimodal Interaction}, 2018, pp. 400--408.

\bibitem{cvxeda2015}
A. Greco, G. Valenza, A. Lanata, E. P. Scilingo, and L. Citi, ``cvxEDA: A Convex Optimization Approach to Electrodermal Activity Processing,'' \textit{IEEE Transactions on Biomedical Engineering}, vol. 63, no. 4, pp. 797--804, 2016.

\end{thebibliography}

\end{document}
